\documentclass[11pt,a4paper]{article}

% --- Kodowanie i język ---
\usepackage[utf8]{inputenc}
\usepackage[T1]{fontenc}
\usepackage[polish]{babel}
\usepackage{lmodern}

% --- Dodatki ---
\usepackage{listings}
\usepackage{xcolor}
\usepackage{geometry}
\usepackage{hyperref}
\geometry{margin=2.3cm}

\hypersetup{
    colorlinks=true,
    linkcolor=blue!50!black,
    urlcolor=blue!50!black
}

% --- Kolory do listingów ---
\definecolor{kolorKomentarz}{rgb}{0.35,0.55,0.35}
\definecolor{kolorSlowo}{rgb}{0.10,0.30,0.70}
\definecolor{kolorString}{rgb}{0.62,0.20,0.18}
\definecolor{kolorTlo}{rgb}{0.975,0.975,0.97}
\definecolor{kolorLinie}{rgb}{0.55,0.55,0.55}

% --- Styl listingu C++ ---
% literate = podmiana znaków, których pdfLaTeX nie umie wyświetlić w kodzie
% (polskie litery oraz myślniki, strzałki itp. z komentarzy w kodzie)
\lstdefinestyle{cpp}{
    language=C++,
    backgroundcolor=\color{kolorTlo},
    basicstyle=\ttfamily\footnotesize,
    keywordstyle=\color{kolorSlowo}\bfseries,
    commentstyle=\color{kolorKomentarz}\itshape,
    stringstyle=\color{kolorString},
    numbers=left,
    numberstyle=\tiny\color{kolorLinie},
    numbersep=8pt,
    showstringspaces=false,
    breaklines=true,
    breakatwhitespace=false,
    frame=single,
    rulecolor=\color{kolorLinie},
    framesep=5pt,
    tabsize=4,
    columns=fullflexible,
    keepspaces=true,
    extendedchars=true,
    literate=%
        {ą}{{\k{a}}}1 {ć}{{\'c}}1 {ę}{{\k{e}}}1 {ł}{{\l}}1 {ń}{{\'n}}1
        {ó}{{\'o}}1 {ś}{{\'s}}1 {ź}{{\'z}}1 {ż}{{\.z}}1
        {Ą}{{\k{A}}}1 {Ć}{{\'C}}1 {Ę}{{\k{E}}}1 {Ł}{{\L}}1 {Ń}{{\'N}}1
        {Ó}{{\'O}}1 {Ś}{{\'S}}1 {Ź}{{\'Z}}1 {Ż}{{\.Z}}1
        {—}{{---}}1 {–}{{--}}1 {„}{{,,}}1 {”}{{''}}1 {’}{{'}}1
        {→}{{$\rightarrow$}}1 {≤}{{$\leq$}}1 {≥}{{$\geq$}}1
        {·}{{$\cdot$}}1 {…}{{...}}1
}
\lstset{style=cpp}

% --- Strona tytułowa ---
\title{\textbf{Sprawozdanie z projektu CML} \\[4pt]
\large Przegląd najciekawszych fragmentów kodu}
\author{Aplikacja do zarządzania biblioteką multimediów (Qt + PostgreSQL)}
\date{\today}

\begin{document}

\maketitle
\tableofcontents
\newpage

% =====================================================================
\section{Wstęp}
% =====================================================================

CML (\emph{Collection Media Library}) to program napisany w C++ z użyciem
biblioteki \textbf{Qt}, który pozwala prowadzić własną bibliotekę multimediów
(książki, gry, seriale, filmy itd.). Dla każdej pozycji można śledzić postęp,
zapisywać sesje (np. ile czasu się spędziło), zaczynać kolejne podejścia do
tego samego tytułu, przypisywać tagi, kategorie i platformy, a na końcu oglądać
rozbudowane statystyki. Dane trzymane są w bazie \textbf{PostgreSQL}.

Kod jest podzielony na trzy warstwy:

\begin{itemize}
    \item \textbf{Warstwa bazy danych} (\texttt{databasemanager.cpp}) --- rozmawia
    bezpośrednio z bazą, wykonuje zapytania SQL i zwraca gotowe dane.
    \item \textbf{Kontroler} (\texttt{appcontroller.cpp}) --- pośrednik między
    bazą a interfejsem. Po każdej udanej zmianie w bazie wysyła sygnał, dzięki
    czemu interfejs sam się odświeża.
    \item \textbf{Interfejs} (pliki w katalogu \texttt{ui/}) --- okna i widżety,
    które widzi użytkownik.
\end{itemize}

W tym sprawozdaniu z każdego pliku \texttt{.cpp} wyciągnęliśmy najciekawsze
fragmenty i opisaliśmy prostym językiem, jak działają i do czego służą. Pominięto
pliki generowane automatycznie przez Qt/CMake (katalog \texttt{build/}, pliki
\texttt{moc\_*}, \texttt{qrc\_*}), bo nie są pisane ręcznie.

% =====================================================================
\section{\texttt{main.cpp} --- punkt startowy programu}
% =====================================================================

Ten plik jest banalnie krótki --- to tylko start aplikacji. Mimo to warto
zrozumieć, co się tutaj dzieje, bo to fundament każdego programu w Qt.

\begin{lstlisting}[caption={Uruchomienie aplikacji}]
int main(int argc, char *argv[])
{
    QApplication a(argc, argv);
    MainWindow w;
    w.show();
    return QCoreApplication::exec();
}
\end{lstlisting}

\textbf{Jak to działa:} \texttt{QApplication} to „silnik” całej aplikacji Qt ---
ustawia obsługę zdarzeń (kliknięcia, naciśnięcia klawiszy itd.). Potem tworzymy
główne okno \texttt{MainWindow} i pokazujemy je przez \texttt{w.show()}.
Najważniejsza jest ostatnia linia: \texttt{exec()} uruchamia tzw. \emph{pętlę
zdarzeń}. To ona sprawia, że program nie kończy się od razu, tylko „stoi" i czeka
na reakcje użytkownika --- dopóki nie zamkniemy okna, program cały czas nasłuchuje,
co robi użytkownik, i odpowiada na to.

% =====================================================================
\section{\texttt{appcontroller.cpp} --- warstwa pośrednia}
% =====================================================================

Kontroler to „kelner” między interfejsem a bazą. Sam prawie nic nie liczy ---
przekazuje żądania do bazy. Ma jednak dwie ważne sztuczki.

\subsection{Sygnał \texttt{daneZmienione()} --- automatyczne odświeżanie}

\begin{lstlisting}[caption={Po udanej zmianie w bazie wysyłamy sygnał}]
int AppController::dodajNoweMedium(const QString &tytul, /* ... */) {
    int id = menedzerBazy.dodajNoweMedium(tytul, /* ... */);
    if (id > 0) {
        emit daneZmienione();   // powiadom interfejs, że coś się zmieniło
    }
    return id;
}
\end{lstlisting}

\textbf{Jak to działa:} prawie każda metoda kontrolera (dodawanie, usuwanie,
edycja) po sukcesie woła \texttt{emit daneZmienione()}. To wysłanie sygnału ---
swego rodzaju ogłoszenie „dane się zmieniły!”. Interfejs jest podłączony pod ten
sygnał (patrz \texttt{mainwindow.cpp}) i automatycznie odświeża drzewo oraz pulpit.

\textbf{Do czego służy:} dzięki temu nie trzeba w 20 miejscach ręcznie pamiętać
o odświeżeniu widoku. Jedno ogłoszenie i całe UI samo się aktualizuje.

\subsection{\texttt{pobierzDaneDlaWykresu} --- zamiana sekund na godziny}

\begin{lstlisting}[caption={Jedyne miejsce z prawdziwą logiką liczenia}]
QList<ActivityStatistic> AppController::pobierzDaneDlaWykresu(int zakres,
                                                  const QString& metryka) {
    auto dane = menedzerBazy.pobierzSuroweDaneStatystyk(zakres, metryka);
    if (metryka == "czas") {
        for (auto& s : dane) s.wartosc /= 3600.0; // sekundy -> godziny
    }
    return dane;
}
\end{lstlisting}

\textbf{Jak to działa:} baza przechowuje czas w sekundach (wygodnie do liczenia),
ale wykres ma pokazywać godziny (czytelniej dla człowieka). Dlatego zanim dane
trafią do wykresu, dzielimy każdą wartość przez 3600 ($3600 = 60 \times 60$ sekund).
To przeliczenie robimy tylko dla metryki „czas” --- przy liczbie sesji czy
zdobytych jednostek nie ma czego przeliczać.

% =====================================================================
\section{\texttt{databasemanager.cpp} --- serce komunikacji z bazą}
% =====================================================================

To najbogatszy plik w projekcie. Tu mieszka cały SQL. Wybraliśmy cztery
najciekawsze fragmenty.

\subsection{\texttt{pobierzWszystkieMultimedia} --- jedno medium = jego najnowsze podejście}

\begin{lstlisting}[caption={Sprytne zapytanie z funkcją okna i DISTINCT ON}]
QSqlQuery query(
    "SELECT DISTINCT ON (m.id) "
    "m.id, m.tytul, p.status, k.id, k.jednostka, "
    "p.wartosc_aktualna, p.wartosc_docelowa, m.data_dodania, "
    "m.id_platformy, p.ocena, m.czy_ulubione, "
    "COALESCE(MAX(COALESCE(d.czas_zakonczenia, d.czas_rozpoczecia)) "
    "         OVER (PARTITION BY p.id), m.data_dodania) AS ostatnia_aktywnosc, "
    "p.numer_podejscia, m.rok_wydania, m.tworcy "
    "FROM multimedia m "
    "LEFT JOIN kategorie k ON m.id_kategorii = k.id "
    "LEFT JOIN podejscia p ON m.id = p.id_medium "
    "LEFT JOIN dziennik_aktywnosci d ON p.id = d.id_podejscia "
    "ORDER BY m.id, p.numer_podejscia DESC;",
    db
);

while (query.next()) {
    auto medium = std::make_shared<Multimedia>();
    medium->id    = query.value(0).toInt();
    medium->tytul = query.value(1).toString();
    // ... (przepisanie pozostałych pól)
    lista.append(medium);
}
\end{lstlisting}

\textbf{Jak to działa --- po kawałku:}
\begin{itemize}
    \item \texttt{DISTINCT ON (m.id)} --- jedno medium może mieć wiele podejść,
    a my chcemy w wyniku tylko jeden wiersz na medium. To polecenie z grupy
    wierszy o tym samym \texttt{m.id} zostawia tylko \emph{pierwszy}. A który jest
    „pierwszy”? Decyduje o tym \texttt{ORDER BY ... numer\_podejscia DESC} ---
    sortujemy podejścia od najnowszego, więc zostaje najnowsze podejście.
    \item \texttt{LEFT JOIN} --- doklejamy dane z innych tabel (kategoria,
    podejście, sesje). „LEFT” gwarantuje, że medium pokaże się \emph{nawet}, jeśli
    nie ma jeszcze kategorii czy żadnej sesji.
    \item \texttt{COALESCE(a, b)} --- zwraca pierwszą wartość, która nie jest pusta
    (NULL). Tutaj: jeśli nie ma daty sesji, bierzemy datę dodania medium.
    \item \texttt{MAX(...) OVER (PARTITION BY p.id)} --- to tzw. funkcja okna.
    Liczy maksimum (najświeższą datę sesji) \emph{w obrębie jednego podejścia},
    ale bez zwijania wierszy do jednego. Dzięki temu obok każdego medium mamy datę
    jego ostatniej aktywności.
\end{itemize}

\textbf{Do czego służy:} to główne zapytanie aplikacji --- jednym strzałem
pobiera listę wszystkich pozycji z najważniejszymi danymi, gotową do wyświetlenia.

\subsection{\texttt{dodajNoweMedium} --- transakcja „wszystko albo nic”}

\begin{lstlisting}[caption={Dwa wpisy do bazy w jednej transakcji}]
db.transaction();   // start transakcji

QSqlQuery query1(db);
// RETURNING id -> PostgreSQL od razu zwróci id nowego wiersza
query1.prepare("INSERT INTO multimedia (tytul, id_kategorii, /* ... */) "
               "VALUES (:tytul, :idKat, /* ... */) RETURNING id");
query1.bindValue(":tytul", tytul);
// id=0 znaczy "nie wybrano" -> zamiast 0 wstawiamy NULL
query1.bindValue(":idKat", idKat > 0 ? QVariant(idKat)
                                     : QVariant(QMetaType::fromType<int>()));

if (!query1.exec() || !query1.next()) {
    db.rollback();   // coś padło -> cofnij WSZYSTKO
    return -1;
}
int noweId = query1.value(0).toInt();   // odczyt id z RETURNING

QSqlQuery query2(db);
query2.prepare("INSERT INTO podejscia (id_medium, numer_podejscia, /* ... */) "
               "VALUES (:id, 1, 0, :cel, 'Planowane')");
query2.bindValue(":id", noweId);
query2.bindValue(":cel", cel);
if (!query2.exec()) { db.rollback(); return -1; }

db.commit();   // dopiero teraz zapis na stałe
return noweId;
\end{lstlisting}

\textbf{Jak to działa:} dodanie medium to tak naprawdę dwa wpisy --- sam tytuł
oraz jego pierwsze podejście. Te dwie operacje muszą się udać razem. Dlatego
opakowujemy je w \emph{transakcję}: \texttt{db.transaction()} ją otwiera,
a \texttt{db.commit()} na końcu zapisuje na stałe. Gdyby cokolwiek po drodze padło,
\texttt{db.rollback()} cofa wszystko --- baza nie zostaje w połowicznym, zepsutym
stanie (np. medium bez podejścia).

\textbf{Dwie dodatkowe sztuczki:} \texttt{RETURNING id} sprawia, że po wstawieniu
od razu dostajemy nadane id (nie trzeba osobnego zapytania). A wstawianie
\texttt{NULL} zamiast zera (gdy użytkownik nic nie wybrał) odróżnia „brak wyboru”
od prawdziwej wartości 0.

\subsection{\texttt{pobierzSuroweDaneStatystyk} --- budowanie zapytania jak z klocków}

\begin{lstlisting}[caption={Część zapytania zależy od wyboru użytkownika}]
QString funkcjaAgregujaca;
if      (metryka == "czas")      funkcjaAgregujaca = "SUM(d.czas_trwania_sekundy)";
else if (metryka == "sesje")     funkcjaAgregujaca = "COUNT(d.id)";
else if (metryka == "jednostki") funkcjaAgregujaca = "SUM(d.przyrost_jednostek)";

// Format daty na osi X dobieramy do długości zakresu
QString formatDaty;
if      (zakresDni <= 30)   formatDaty = "'DD.MM'";    // pojedyncze dni
else if (zakresDni <= 365)  formatDaty = "'MM.YYYY'";  // miesiące
else if (zakresDni > 1000)  formatDaty = "'YYYY'";     // lata
else                        formatDaty = "'YYYY-MM'";

// %1, %2, %3 to FRAGMENTY SQL, dlatego nie da się ich podać przez bindValue
QString sql = QString(R"(
    SELECT TO_CHAR(/* data */, %1) as kategoria_czasu,
           m.tytul, %2 as wartosc
    FROM dziennik_aktywnosci d
    JOIN podejscia p ON d.id_podejscia = p.id
    JOIN multimedia m ON p.id_medium = m.id
    WHERE /* data */ >= CURRENT_DATE - INTERVAL '%3 days'
    GROUP BY kategoria_czasu, m.tytul
)").arg(formatDaty).arg(funkcjaAgregujaca).arg(zakresDni);
\end{lstlisting}

\textbf{Jak to działa:} wykres aktywności może pokazywać różne rzeczy (czas, liczbę
sesji, zdobyte jednostki) i różne zakresy (tydzień, rok, cała historia). Zamiast
pisać osobne zapytanie na każdą kombinację, sklejamy zapytanie z gotowych
fragmentów --- najpierw wybieramy odpowiednią funkcję sumującą i format daty,
a potem wstawiamy je w szablon przez \texttt{QString::arg} (\texttt{\%1},
\texttt{\%2}, \texttt{\%3}).

\textbf{Ważny szczegół:} normalnie wartości wstawia się bezpiecznie przez
\texttt{bindValue}, ale tutaj wstawiamy \emph{fragmenty samego SQL} (nazwę funkcji,
wzór daty), a tego \texttt{bindValue} nie potrafi --- ono podstawia tylko wartości,
nie kawałki polecenia. Dlatego tu sklejamy tekst ręcznie.

\subsection{\texttt{pobierzCiekawostkiStatystyczne} --- zabawne statystyki}

\begin{lstlisting}[caption={Pora doby i dominująca dekada}]
// Pora doby: przypisz każdej sesji etykietę wg godziny, weź najpopularniejszą
q8.prepare(R"(
    SELECT pora FROM (
        SELECT CASE
            WHEN EXTRACT(HOUR FROM czas_rozpoczecia) BETWEEN 6 AND 11  THEN 'Rano'
            WHEN EXTRACT(HOUR FROM czas_rozpoczecia) BETWEEN 12 AND 17 THEN 'Popoludnie'
            WHEN EXTRACT(HOUR FROM czas_rozpoczecia) BETWEEN 18 AND 23 THEN 'Wieczor'
            ELSE 'Noc'
        END AS pora, SUM(czas_trwania_sekundy) AS suma
        FROM dziennik_aktywnosci
        GROUP BY pora ORDER BY suma DESC LIMIT 1
    ) AS najaktywniejsza
)");

// Dominujaca dekada: 1997 / 10 * 10 = 1990 (dzielenie calkowite ucina koncowke)
q6.prepare(R"(
    SELECT (rok_wydania / 10) * 10 AS dekada, COUNT(*) AS ilosc
    FROM multimedia WHERE rok_wydania > 0
    GROUP BY dekada ORDER BY ilosc DESC LIMIT 1
)");
\end{lstlisting}

\textbf{Jak to działa:} ta funkcja zbiera „ciekawostki” do panelu statystyk ---
każda to małe, niezależne zapytanie. W pierwszym przykładzie \texttt{CASE ... END}
działa jak ciąg „jeżeli... to...” --- patrzy na godzinę rozpoczęcia sesji
(\texttt{EXTRACT(HOUR ...)}) i przykleja etykietę (Rano, Popołudnie, Wieczór, Noc).
Potem sumujemy czas dla każdej pory i bierzemy tę z największą sumą --- stąd wiemy,
czy użytkownik to „Nocny Marek”, czy „Ranny Ptaszek”.

W drugim przykładzie liczymy najpopularniejszą dekadę. Sztuczka to dzielenie
całkowite: \texttt{1997 / 10} daje 199 (końcówka odpada), a po pomnożeniu przez 10
wychodzi 1990 --- czyli „lata 90.”.

\textbf{Ciekawostka:} w tej funkcji jest też zapytanie z \texttt{EXTRACT(ISODOW ...)},
które zwraca numer dnia tygodnia (1 = poniedziałek, 7 = niedziela) --- aplikacja
podaje na tej podstawie najbardziej aktywny dzień tygodnia.

% =====================================================================
\section{\texttt{mainwindow.cpp} --- okno główne i centrum dowodzenia}
% =====================================================================

Główne okno spina wszystkie panele w jedną całość i decyduje, który widok jest
aktualnie pokazywany.

\subsection{Łączenie sygnałów lambdami}

\begin{lstlisting}[caption={Reakcje na zdarzenia z paneli}]
// Klik w drzewie nawigacji -> pokaz szczegoly wybranego medium
connect(panelNawigacji, &NavigationPanelWidget::zadaniePokazaniaSzczegolow,
        this, [this](int id) { pokazSzczegolyMedium(id); });

// Po KAZDEJ zmianie danych w bazie odswiez drzewo i pulpit
connect(&appController, &AppController::daneZmienione, this, [this]() {
    panelNawigacji->odswiezDrzewo();
    pulpitWidget->odswiezStatystykiGlowne();
});
\end{lstlisting}

\textbf{Jak to działa:} \texttt{connect} łączy sygnał (coś się stało) ze slotem
(co zrobić). Tutaj używamy \emph{lambd} --- małych funkcji pisanych w miejscu.
Pierwsze połączenie mówi: „gdy panel poprosi o pokazanie szczegółów medium,
wywołaj \texttt{pokazSzczegolyMedium}”. Drugie to wspomniana wcześniej magia
odświeżania --- okno nasłuchuje sygnału \texttt{daneZmienione()} z kontrolera
i samo aktualizuje widok.

\textbf{Do czego służy:} to spina całą aplikację. Panele nie muszą się nawzajem
znać --- wysyłają sygnały, a okno główne decyduje, co z nimi zrobić.

\subsection{Szablon \texttt{otworzDialogZarzadzania} --- jedna metoda na wiele okien}

\begin{lstlisting}[caption={Generyczne otwieranie dialogów}]
template <typename Dialog>
void MainWindow::otworzDialogZarzadzania() {
    Dialog dialog(appController, this);
    dialog.exec();                              // pokaz i czekaj az uzytkownik zamknie
    pulpitWidget->odswiezStatystykiGlowne();    // po zamknieciu odswiez pulpit
}

// Uzycie - ten sam kod dla trzech roznych okien:
connect(ui->actionKategorie, &QAction::triggered, this,
        [this]() { otworzDialogZarzadzania<CategoriesDialog>(); });
connect(ui->actionPlatformy, &QAction::triggered, this,
        [this]() { otworzDialogZarzadzania<PlatformsDialog>(); });
connect(ui->actionTagi, &QAction::triggered, this,
        [this]() { otworzDialogZarzadzania<TagsDialog>(); });
\end{lstlisting}

\textbf{Jak to działa:} dialogi Kategorie, Platformy i Tagi otwiera się tak samo
(utwórz, pokaż, po zamknięciu odśwież pulpit). Zamiast pisać trzy prawie
identyczne metody, używamy \emph{szablonu} --- \texttt{template <typename Dialog>}
oznacza „wstaw tu dowolny typ okna”. Przy wywołaniu podajemy konkretny typ
w nawiasach ostrych, np. \texttt{<CategoriesDialog>}.

\textbf{Do czego służy:} mniej powtórzonego kodu. Jeden szablon obsługuje wszystkie
trzy okna --- a jak dojdzie czwarte, wystarczy jedna linijka.

% =====================================================================
\section{\texttt{dashboardwidget.cpp} --- pulpit startowy}
% =====================================================================

Pulpit to pierwszy ekran po uruchomieniu: wykres kołowy ze statusami biblioteki
oraz kafelki ostatnio używanych pozycji.

\subsection{Wykres kołowy statusów}

\begin{lstlisting}[caption={Budowa wykresu kołowego}]
QPieSeries *seria = new QPieSeries();
seria->append(Status::Planowane, statystyki.value(Status::Planowane, 0));
seria->append(Status::WTrakcie,  statystyki.value(Status::WTrakcie, 0));
// ... pozostale statusy

// Kolory wycinkow wg kolejnosci dodania
seria->slices().at(0)->setColor(QColor(0x7f, 0x8c, 0x8d));
seria->slices().at(3)->setColor(QColor(0x27, 0xae, 0x60));

// Doklej liczbe do etykiety: "Planowane" -> "Planowane (15)"
for (auto wycinek : seria->slices()) {
    wycinek->setLabel(QString("%1 (%2)").arg(wycinek->label()).arg(wycinek->value()));
}

QChart *wykres = new QChart();
wykres->addSeries(seria);
wykres->setTitle("Status Biblioteki");
ui->wykresOwalny->setChart(wykres);
\end{lstlisting}

\textbf{Jak to działa:} \texttt{QPieSeries} to „ciasto”, a każde \texttt{append} dokłada
wycinek --- nazwę statusu i liczbę pozycji w tym statusie. Potem ręcznie ustawiamy
kolory (np. zielony dla ukończonych) i w pętli dopisujemy do każdej etykiety
liczbę w nawiasie. Na końcu wkładamy serię do wykresu i pokazujemy.
Na samej górze funkcji jest jeszcze zabezpieczenie: jeśli biblioteka jest pusta,
funkcja kończy się od razu --- nie ma czego rysować.

\subsection{Kafelki „ostatnio aktywne” z łamaniem wierszy}

\begin{lstlisting}[caption={Czyszczenie siatki i układanie kafelków po 2 w rzędzie}]
// Najpierw usun stare kafelki, zeby nie dokladac do poprzednich
QLayoutItem *element;
while ((element = ui->gridOstatnie->takeAt(0)) != nullptr) {
    delete element->widget();
    delete element;
}

const QList<int> ostatnieId = appController.pobierzOstatnioAktywne(6);
int wiersz = 0, kolumna = 0;
for (int id : ostatnieId) {
    // ... (znajdz medium o tym id) ...
    QPushButton *btnKafel = new QPushButton(m->tytul, this);
    connect(btnKafel, &QPushButton::clicked, this, [this, id]() {
        emit zadaniePokazaniaSzczegolow(id);
    });
    ui->gridOstatnie->addWidget(btnKafel, wiersz, kolumna);

    kolumna++;
    if (kolumna > 1) { kolumna = 0; wiersz++; }   // maks. 2 kafelki w rzedzie
}
\end{lstlisting}

\textbf{Jak to działa:} najpierw pętla \texttt{takeAt(0)} wyjmuje i kasuje stare
kafelki --- bez tego przy każdym odświeżeniu dokładałyby się nowe na stare.
Potem dla sześciu ostatnio używanych pozycji tworzymy przyciski. Łamanie wierszy
to prosta arytmetyka: zwiększamy numer kolumny, a gdy przekroczy 1 (czyli mamy
już 2 kafelki w rzędzie), wracamy do kolumny 0 i schodzimy do następnego wiersza.

% =====================================================================
\section{\texttt{detailswidget.cpp} --- ekran szczegółów pozycji}
% =====================================================================

To jeden z najbardziej rozbudowanych widoków --- pokazuje wszystko o jednej
pozycji: postęp, historię podejść i sesji, ocenę, recenzję. Wybraliśmy cztery
fragmenty.

\subsection{Dialog sesji ze stoperem}

\begin{lstlisting}[caption={Stoper zsynchronizowany ze spinnerami h/m/s}]
int lacznieSekund = sekundy;

// Lambda: jednoczesnie odswieza wyswietlacz stopera i trzy spinnery
auto odswiezWyswietlacz = [&]() {
    int h = lacznieSekund / 3600;
    int m = (lacznieSekund % 3600) / 60;
    int s = lacznieSekund % 60;

    // blockSignals: programowa zmiana spinnera NIE wywola naszego kodu z powrotem
    // (inaczej powstalaby nieskonczona petla)
    spinH->blockSignals(true); spinM->blockSignals(true); spinS->blockSignals(true);
    spinH->setValue(h); spinM->setValue(m); spinS->setValue(s);
    spinH->blockSignals(false); spinM->blockSignals(false); spinS->blockSignals(false);

    lblStoper->setText(QString("%1:%2:%3")
        .arg(h,2,10,QChar('0')).arg(m,2,10,QChar('0')).arg(s,2,10,QChar('0')));
};

// Timer tykajacy co sekunde -> dodaje 1 sekunde i odswieza widok
QTimer* czasomierz = new QTimer(&dialog);
connect(czasomierz, &QTimer::timeout, [&]() { lacznieSekund++; odswiezWyswietlacz(); });

// Jeden przycisk pelni dwie role: Start i Pauza
connect(btnStart, &QPushButton::clicked, [&]() {
    if (czasomierz->isActive()) czasomierz->stop();   // bylo wlaczone -> pauza
    else                        czasomierz->start(1000); // bylo wylaczone -> start
});
\end{lstlisting}

\textbf{Jak to działa:} okno dodawania sesji ma stoper, który tyka co sekundę
(\texttt{QTimer} z odstępem 1000\,ms), oraz trzy pola na ręczne wpisanie godzin,
minut i sekund. Obie rzeczy są ze sobą zsynchronizowane --- gdy stoper tyka,
pola się zmieniają, a gdy ręcznie zmienisz pola, zmienia się stoper.

\textbf{Najważniejsza sztuczka --- \texttt{blockSignals}:} gdyby program ustawiał
wartość pola, a to pole automatycznie wywoływało kod przeliczający czas, ten kod
znowu ustawiłby pole i tak w kółko (nieskończona pętla). \texttt{blockSignals(true)}
na chwilę „wycisza” pola, żeby zmiana programowa nie odpaliła reakcji.

\subsection{Drzewo historii --- podejścia i sesje}

\begin{lstlisting}[caption={Podejście jako rodzic, sesja jako dziecko}]
for (const auto& p : historia) {
    // Podejscie = wezel-rodzic, dodawany do korzenia drzewa
    QTreeWidgetItem *wezelPodejscia = new QTreeWidgetItem(ui->treeHistoria);
    wezelPodejscia->setText(0, QString("Podejscie #%1 (%2)").arg(p.numer).arg(p.status));
    // Ukryte dane: TYP wezla i ID rekordu w bazie (czytane po kliknieciu)
    wezelPodejscia->setData(0, ROLA_TYP, static_cast<int>(TypHistoriiElementu::Podejscie));
    wezelPodejscia->setData(0, ROLA_ID, p.id);

    for (int i = 0; i < p.sesje.size(); ++i) {
        const auto& s = p.sesje[i];
        // Sesja = wezel-dziecko podpiety pod podejscie
        QTreeWidgetItem *wezelSesji = new QTreeWidgetItem(wezelPodejscia);
        wezelSesji->setText(0, QString("Sesja #%1").arg(p.sesje.size() - i));
        wezelSesji->setData(0, ROLA_ID, s.id);
        // HTML ze szczegolami - pokazywany w panelu na dole po kliknieciu
        wezelSesji->setData(0, Qt::UserRole, infoSesja);
    }
}
ui->treeHistoria->expandAll();
\end{lstlisting}

\textbf{Jak to działa:} historia ma dwa poziomy: podejścia (rodzice) i należące
do nich sesje (dzieci). Tworzymy więc drzewo --- nowy \texttt{QTreeWidgetItem}
z drzewem jako rodzicem to węzeł najwyższego poziomu, a z innym węzłem jako
rodzicem --- dziecko.

\textbf{Sprytny trik --- ukryte dane:} w każdym węźle chowamy dodatkowe informacje
przez \texttt{setData} z własnymi „rolami” (\texttt{ROLA\_TYP}, \texttt{ROLA\_ID}).
Użytkownik ich nie widzi, ale gdy kliknie węzeł, program odczytuje te dane i wie,
czy kliknięto podejście czy sesję i którego rekordu w bazie dotyczy. Pod rolą
\texttt{Qt::UserRole} chowamy też gotowy opis HTML do pokazania w panelu na dole.

\subsection{Dynamiczne pokazywanie i ukrywanie kontrolek}

\begin{lstlisting}[caption={Inny układ dla pozycji zakończonej i aktywnej}]
bool czyZakonczone = (medium->status == Status::Ukonczone ||
                      medium->status == Status::Porzucone);
if (czyZakonczone) {
    ui->lblWielkiStatus->setText(medium->status.toUpper());
    ui->frameKontrolkiPostepu->hide();   // brak edycji postepu
    ui->comboDetaleStatus->hide();
    ui->btnNowePodejscie->show();         // za to: "Nowe podejscie" / "Cofnij"
} else {
    ui->lblWielkiStatus->hide();
    ui->frameKontrolkiPostepu->show();   // pelna edycja postepu
    ui->comboDetaleStatus->show();
}

// Licznik "Brak postepu od: X dni" - po 30 dniach robi sie czerwony
int dniBezAkcji = odKiedyBrakPostepu.daysTo(teraz);
QString kolor = (dniBezAkcji > 30) ? "#e74c3c"
                                   : palette().color(QPalette::WindowText).name();
\end{lstlisting}

\textbf{Jak to działa:} ten sam ekran wygląda inaczej w zależności od stanu
pozycji. Jeśli tytuł jest ukończony lub porzucony, chowamy kontrolki postępu
(nie ma już czego zmieniać) i pokazujemy duży napis ze statusem oraz przyciski
„Nowe podejście” i „Cofnij zakończenie”. Dla aktywnego tytułu jest odwrotnie ---
pełna edycja postępu.

\textbf{Miły detal:} aplikacja liczy, ile dni minęło bez postępu. Jeśli ponad 30,
licznik robi się czerwony (\texttt{\#e74c3c}) --- delikatne przypomnienie, że
o tytule zapomniałeś.

\subsection{Automatyczne ukończenie po osiągnięciu celu}

\begin{lstlisting}[caption={Pytanie o automatyczne zamknięcie pozycji}]
// Jesli aktualna == docelowa, a uzytkownik sam nie ustawil "Ukonczone",
// pytamy, czy zamknac pozycje automatycznie.
if (nowyStatus != Status::Ukonczone &&
    appController.czyOsiagnietoCel(nowaAktualna, nowaDocelowa)) {
    auto odp = QMessageBox::question(this, "Automatyczne ukonczenie",
        "Wartosc aktualna zrownala sie z docelowa. "
        "Po zapisaniu status zmieni sie na 'Ukonczone'. Kontynuowac?",
        QMessageBox::Yes | QMessageBox::No);
    if (odp == QMessageBox::No) return;
    nowyStatus = Status::Ukonczone;
}

if (appController.aktualizujPostep(aktualneIdMedium, nowyStatus,
                                  nowaAktualna, nowaDocelowa, ocena)) {
    ustawMedium(aktualneIdMedium);   // odswiez widok
    emit daneZaktualizowane();       // powiadom reszte aplikacji
}
\end{lstlisting}

\textbf{Jak to działa:} przy zapisie program sprawdza, czy wartość aktualna
zrównała się z docelową (np. przeczytałeś 350 z 350 stron). Jeśli tak, a status
nie jest jeszcze „Ukończone”, wyskakuje pytanie, czy automatycznie zamknąć pozycję.
To wygoda --- użytkownik nie musi pamiętać o ręcznej zmianie statusu, ale ma też
prawo odmówić (\texttt{return} przerywa zapis).

% =====================================================================
\section{\texttt{navigationpanelwidget.cpp} --- panel z drzewem po lewej}
% =====================================================================

To panel boczny z drzewem pozycji oraz wyszukiwarką. Pozwala grupować bibliotekę
na kilka sposobów.

\subsection{Grupowanie drzewa na 6 sposobów}

\begin{lstlisting}[caption={Jedno medium trafia do odpowiedniego folderu}]
const int trybGrupowania = ui->comboGrupowanie->currentIndex();
// 0=Kategoria 1=Status 2=Platforma 3=Data 4=Tagi 5=Typ nosnika

for (const auto& medium : listaMultimediow) {
    // Tagi: jedno medium moze trafic do KILKU folderow (relacja wiele-do-wielu)
    if (trybGrupowania == 4) {
        QStringList tagiMedium = mapaTagow.value(medium->id);
        if (tagiMedium.isEmpty()) tagiMedium << "Brak tagow";
        for (const QString& nazwaGrupy : tagiMedium) {
            // ... dodaj medium do folderu o nazwie nazwaGrupy ...
        }
        continue;
    }

    // Pozostale tryby: jedno medium = jeden folder
    QString nazwaGrupy;
    switch (trybGrupowania) {
    case 0: nazwaGrupy = slownikKategorii.value(medium->idKategorii, "Brak kategorii"); break;
    case 1: nazwaGrupy = medium->status; break;
    case 3: {
        QLocale polski(QLocale::Polish, QLocale::Poland);
        QDate data = medium->dataDodania.date();
        nazwaGrupy = QString("%1 %2").arg(polski.standaloneMonthName(data.month()))
                                     .arg(data.year()).toUpper();
        break;
    }
    // ...
    }
    // ... znajdz lub utworz folder i wstaw do niego medium ...
}
\end{lstlisting}

\textbf{Jak to działa:} użytkownik wybiera z listy rozwijanej sposób grupowania,
a jego numer steruje całą funkcją. Konstrukcja \texttt{switch} dla każdego trybu
ustala, do jakiej grupy (folderu) trafia dane medium --- po kategorii, statusie,
platformie, miesiącu dodania itd.

\textbf{Wyjątek --- tagi:} tagi to relacja „wiele-do-wielu” (jedno medium może mieć
kilka tagów), więc obsługujemy je osobno --- to samo medium dodajemy do każdego
folderu z jego tagiem. Dlatego ten przypadek ma własną pętlę i \texttt{continue}.

\subsection{Wyszukiwarka filtrująca „w miejscu”}

\begin{lstlisting}[caption={Ukrywanie niepasujących pozycji bez przebudowy drzewa}]
for (int i = 0; i < ui->kategorie->topLevelItemCount(); ++i) {
    QTreeWidgetItem *kategoria = ui->kategorie->topLevelItem(i);   // folder
    bool maPasujaceElementy = false;

    for (int j = 0; j < kategoria->childCount(); ++j) {
        QTreeWidgetItem *dziecko = kategoria->child(j);            // medium
        const bool pasuje = dziecko->text(0).contains(tekst, Qt::CaseInsensitive);
        dziecko->setHidden(!pasuje);
        if (pasuje) maPasujaceElementy = true;
    }

    if (tekst.isEmpty()) {
        kategoria->setHidden(false);                // pusty tekst -> pokaz wszystko
    } else {
        kategoria->setHidden(!maPasujaceElementy);  // folder bez trafien znika
        if (maPasujaceElementy) kategoria->setExpanded(true);
    }
}
\end{lstlisting}

\textbf{Jak to działa:} zamiast przebudowywać całe drzewo przy każdej wpisanej
literze, po prostu chowamy i pokazujemy istniejące elementy. Przechodzimy po
folderach, a w każdym po pozycjach. \texttt{contains(..., Qt::CaseInsensitive)}
sprawdza, czy tytuł zawiera wpisany tekst (bez znaczenia wielkość liter).
Niepasujące pozycje chowamy. Jeśli w folderze nie zostało żadne trafienie,
chowamy też cały folder; a jeśli coś zostało --- rozwijamy go, żeby od razu było
widać wynik. Pusty tekst przywraca pełny widok.

% =====================================================================
\section{\texttt{multimediaformwidget.cpp} --- formularz dodawania i edycji}
% =====================================================================

Formularz służy do dodawania nowej pozycji i edycji istniejącej. Ten sam
formularz obsługuje oba przypadki.

\subsection{Lista tagów z checkboxami}

\begin{lstlisting}[caption={Budowa listy i odczyt zaznaczonych tagów}]
// Budowa: kazdy tag = pozycja z checkboxem
for (const auto& tag : wszystkieTagi) {
    QListWidgetItem* item = new QListWidgetItem(tag.second, ui->listTagi);
    item->setFlags(item->flags() | Qt::ItemIsUserCheckable);  // wlacz checkbox
    item->setData(Qt::UserRole, tag.first);                   // schowaj ID tagu
    // W trybie edycji juz przypisane tagi sa od razu zaznaczone
    if (czyTrybEdycji && tagiTegoMedium.contains(tag.second))
        item->setCheckState(Qt::Checked);
    else
        item->setCheckState(Qt::Unchecked);
}

// Odczyt przy zapisie: zbieramy ID zaznaczonych tagow
QList<int> wybraneTagi;
for (int i = 0; i < ui->listTagi->count(); ++i)
    if (ui->listTagi->item(i)->checkState() == Qt::Checked)
        wybraneTagi.append(ui->listTagi->item(i)->data(Qt::UserRole).toInt());
appController.ustawTagiDlaMedium(idDoTagow, wybraneTagi);
\end{lstlisting}

\textbf{Jak to działa:} każdy istniejący tag dostaje swój wiersz na liście
z polem do zaznaczenia (\texttt{Qt::ItemIsUserCheckable} włącza checkbox).
Tak jak w drzewie historii, chowamy ID tagu pod \texttt{Qt::UserRole} ---
na liście widać tylko nazwę, ale program pamięta, który to tag.
W trybie edycji tagi już przypisane do pozycji są od razu zaznaczone.
Przy zapisie przechodzimy po liście i zbieramy ID tylko tych zaznaczonych.

\subsection{Etykieta celu zależna od kategorii}

\begin{lstlisting}[caption={Podpis celu zmienia się razem z jednostką kategorii}]
// Pod kazda kategoria w combo chowamy jej jednostke (np. "strony")
ui->comboNowaKategoria->setItemData(idx, jedn, ROLA_JEDNOSTKA);

// Po zmianie kategorii podmieniamy etykiete celu
void MultimediaFormWidget::obsluzZmianeKategorii(int indeks) {
    QString jednostka = ui->comboNowaKategoria->itemData(indeks, ROLA_JEDNOSTKA).toString();
    ui->labJednostka->setText(jednostka.isEmpty() ? "Cel:" : "Cel (" + jednostka + "):");
}
\end{lstlisting}

\textbf{Jak to działa:} każda kategoria ma swoją jednostkę (książki --- strony,
seriale --- odcinki). Przy budowie listy rozwijanej chowamy tę jednostkę pod
kategorią. Gdy użytkownik wybierze inną kategorię, odczytujemy jej jednostkę
i podmieniamy podpis pola celu, np. na „Cel (strony):” albo „Cel (odcinki):”.

\textbf{Do czego służy:} drobiazg, ale formularz od razu mówi, w czym podajesz
cel --- nie trzeba zgadywać, czy chodzi o strony, godziny czy odcinki.

% =====================================================================
\section{\texttt{statisticswidget.cpp} --- panel statystyk i wykresów}
% =====================================================================

Najbardziej „efekciarski” widok --- pełno wykresów i zestawień. Wybraliśmy cztery
fragmenty.

\subsection{Wykres słupkowy z rankingiem wewnątrz słupka}

\begin{lstlisting}[caption={Każdy słupek pokazuje swoich liderów}]
const int TOP_N = 4;   // miejsca 1-4, reszta laduje w "Pozostale"

for (const QString& d : unikalneDaty) {
    // (wartosc, tytul) dla tego slupka, posortowane malejaco
    QList<QPair<double, QString>> wDniu;
    for (auto it = mapaSerii.begin(); it != mapaSerii.end(); ++it) {
        const double val = it.value().value(d, 0.0);
        if (val > 0.0) wDniu.append({val, it.key()});
    }
    std::sort(wDniu.begin(), wDniu.end(), [](const auto& a, const auto& b) {
        return a.first > b.first;   // od najwiekszej wartosci
    });

    // Miejsca 1..TOP_N - po jednym tytule na miejsce (lub 0, gdy brak)
    for (int r = 0; r < TOP_N; ++r) {
        if (r < wDniu.size()) *setySlotow[r] << wDniu[r].first;
        else                  *setySlotow[r] << 0.0;
    }
    // Slot zbiorczy "Pozostale" - suma wszystkiego poza pierwsza czworka
    double sumaReszty = 0.0;
    for (int r = TOP_N; r < wDniu.size(); ++r) sumaReszty += wDniu[r].first;
    *setySlotow[TOP_N] << sumaReszty;
}
\end{lstlisting}

\textbf{Jak to działa:} to wykres słupkowy „skumulowany” --- każdy słupek
(np. jeden dzień) składa się z kilku kolorowych kawałków. Dla każdego słupka
osobno sortujemy tytuły malejąco i przypisujemy je do „miejsc” 1--4, a wszystko
poniżej sumujemy w jeden kawałek „Pozostałe”.

\textbf{Najważniejsze:} kolor oznacza \emph{miejsce w rankingu}, a nie konkretny
tytuł --- pierwsze miejsce zawsze złote, drugie srebrne, trzecie brązowe. Dzięki
temu każdy słupek pokazuje swoich własnych liderów, a po najechaniu myszą
w dymku widać, jaki to tytuł i ile dokładnie.

\subsection{Pływający dymek podążający za kursorem}

\begin{lstlisting}[caption={Filtr zdarzeń przesuwający etykietę i odbijający ją od krawędzi}]
bool StatisticsWidget::eventFilter(QObject *watched, QEvent *event) {
    if (watched == ui->wykresSlupkowyAktywnosci->viewport()
        && event->type() == QEvent::MouseMove) {
        if (!etykietaDymka->isHidden()) {
            QMouseEvent *m = static_cast<QMouseEvent*>(event);
            int x = m->pos().x() + 15;   // 15px obok kursora, by go nie zaslaniac
            int y = m->pos().y() + 15;

            // Jesli dymek wychodzilby za krawedz - przerzuc na druga strone kursora
            auto *vp = ui->wykresSlupkowyAktywnosci->viewport();
            if (x + etykietaDymka->width()  > vp->width())
                x = m->pos().x() - etykietaDymka->width()  - 15;
            if (y + etykietaDymka->height() > vp->height())
                y = m->pos().y() - etykietaDymka->height() - 15;

            etykietaDymka->move(x, y);
            etykietaDymka->raise();   // dymek zawsze na wierzchu
        }
    }
    return QWidget::eventFilter(watched, event);
}
\end{lstlisting}

\textbf{Jak to działa:} \texttt{eventFilter} to „podsłuch” zdarzeń --- przechwytuje
ruchy myszy nad wykresem. Gdy dymek jest widoczny, ustawiamy go 15 pikseli obok
kursora (żeby go nie zasłaniać). Jeśli dymek wyszedłby poza prawą lub dolną
krawędź wykresu, przerzucamy go na drugą stronę kursora --- dzięki temu zawsze
mieści się w obszarze wykresu. \texttt{raise()} pilnuje, żeby dymek był na wierzchu.

\subsection{Wykres radarowy (profil tagów)}

\begin{lstlisting}[caption={Wielokąt z najczęstszych tagów}]
// Radar (wielokat) ma sens dopiero od 3 wierzcholkow
if (nTagow > 2) {
    for (int i = 0; i < nTagow; ++i) {
        // normalizacja do 0-100% wzgledem najwiekszego tagu
        const qreal norm = (qreal(topTagi[i].value("czasSekundy").toLongLong())
                            / qreal(maxCzasTagu)) * 100.0;
        radarSeries->append(i, norm);
        axisAngular->append(topTagi[i].value("tag").toString(), i);
    }
    // Domkniecie figury: powtarzamy pierwszy punkt na koncu (kat = 360 stopni)
    const qreal normZamkniecie = (qreal(topTagi[0].value("czasSekundy").toLongLong())
                                  / qreal(maxCzasTagu)) * 100.0;
    radarSeries->append(nTagow, normZamkniecie);
    axisAngular->setRange(0, nTagow);
} else {
    radarSeries->append(0, 0);
    axisAngular->append("Malo danych", 1);
}
\end{lstlisting}

\textbf{Jak to działa:} wykres radarowy to wielokąt, którego wierzchołki to
najczęstsze tagi, a „wysięg” każdego wierzchołka pokazuje, ile czasu poszło na
dany tag. Wartości \emph{normalizujemy} do skali 0--100\% --- największy tag dostaje
100\%, reszta proporcjonalnie. To pozwala porównać tagi niezależnie od tego, czy
liczymy w minutach, czy godzinach.

\textbf{Dwie pułapki, które kod obsługuje:} po pierwsze, żeby figura była zamknięta,
trzeba powtórzyć pierwszy punkt na końcu (inaczej zostałaby „dziura” między
ostatnim a pierwszym wierzchołkiem). Po drugie, wielokąt ma sens dopiero od 3
wierzchołków --- przy 1--2 tagach pokazujemy zastępcze „Mało danych” zamiast
zdegenerowanej figury.

\subsection{„Kupka Wstydu” --- losowa propozycja}

\begin{lstlisting}[caption={Losowanie pozycji do nadrobienia}]
// disconnect przed connect: ta funkcja biega przy kazdym wejsciu na zakladke.
// Bez tego polaczenia by sie mnozyly i jeden klik wywolywalby lambde wielokrotnie.
disconnect(ui->btnProponujAktywnosc, nullptr, this, nullptr);
connect(ui->btnProponujAktywnosc, &QPushButton::clicked, this,
        [this, wszystkieWKupce]() {
    if (wszystkieWKupce.isEmpty()) {
        QMessageBox::information(this, "Kupka Wstydu", "Brak pozycji do zaproponowania.");
        return;
    }
    int indeks = QRandomGenerator::global()->bounded(wszystkieWKupce.size());
    emit zadaniePokazaniaSzczegolow(wszystkieWKupce.at(indeks)->id);
});
\end{lstlisting}

\textbf{Jak to działa:} „Kupka Wstydu” to lista zaczętych, ale niedokończonych
pozycji. Przycisk „Zaproponuj aktywność” losuje jedną z nich:
\texttt{QRandomGenerator::global()->bounded(n)} zwraca losową liczbę od 0 do
$n-1$, czyli losowy indeks na liście. Potem wysyłamy sygnał, żeby pokazać
szczegóły wylosowanej pozycji.

\textbf{Ważny detal --- \texttt{disconnect}:} ta funkcja uruchamia się za każdym
razem, gdy wchodzisz na zakładkę. Gdybyśmy tylko dodawali połączenie
(\texttt{connect}), to po pięciu wejściach jeden klik wywołałby losowanie pięć
razy. Dlatego najpierw \texttt{disconnect} kasuje stare połączenia, a dopiero
potem robimy świeże.

% =====================================================================
\section{\texttt{categoriesdialog.cpp} --- okno zarządzania kategoriami}
% =====================================================================

Okno do dodawania, edycji i usuwania kategorii. Najciekawsze jest usuwanie ---
bo kategoria może mieć przypisane pozycje.

\subsection{Usuwanie z wyborem losu zawartości}

\begin{lstlisting}[caption={Własne przyciski i rozpoznanie, który kliknięto}]
QMessageBox msgBox(this);
msgBox.setText("Usuwasz kategorie: <b>" + nazwa + "</b><br>Co robimy z zawartoscia?");

// Zapamietujemy wskazniki, by potem rozpoznac KTORY przycisk kliknieto
QPushButton *btnPrzenies     = msgBox.addButton("Przenies do 'Brak'", QMessageBox::ActionRole);
QPushButton *btnUsunWszystko = msgBox.addButton("Usun wszystko",      QMessageBox::DestructiveRole);
msgBox.addButton("Anuluj", QMessageBox::RejectRole);

msgBox.exec();

QAbstractButton *klikniety = msgBox.clickedButton();
if (klikniety != btnPrzenies && klikniety != btnUsunWszystko) return;  // anulowano

const bool usunZDetalami = (klikniety == btnUsunWszystko);
if (appController.usunKategorie(idKat, usunZDetalami)) {
    wypelnijTabele();
}
\end{lstlisting}

\textbf{Jak to działa:} zwykłe okienko ma przyciski Tak/Nie, ale tutaj potrzebne
są trzy opcje, więc dodajemy własne przyciski przez \texttt{addButton}.
Zapamiętujemy wskaźniki do nich, bo po zamknięciu \texttt{clickedButton()} mówi
tylko, który obiekt kliknięto --- porównujemy go z naszymi wskaźnikami, żeby wiedzieć,
co użytkownik wybrał. „Przenieś do Brak” odpina pozycje od kategorii (zostają
w bibliotece), a „Usuń wszystko” kasuje je razem z kategorią.

\subsection{Edytowalna lista jednostek}

\begin{lstlisting}[caption={Podpowiedzi z bazy, ale można wpisać własną}]
QComboBox comboJednostka(&dialog);
comboJednostka.setEditable(true);                                  // mozna wpisac wlasna
comboJednostka.addItems(appController.pobierzUnikalneJednostki()); // + podpowiedzi z bazy
// ...
const QString jednostka = comboJednostka.currentText().trimmed();
// Gdy puste -> domyslnie "szt."
appController.dodajKategorie(nazwa, jednostka.isEmpty() ? "szt." : jednostka);
\end{lstlisting}

\textbf{Jak to działa:} lista jednostek jest \emph{edytowalna}
(\texttt{setEditable(true)}) --- użytkownik może wybrać gotową podpowiedź
(pobraną z bazy z już istniejących kategorii) albo wpisać zupełnie nową, np.
„rozdziały”. Gdy zostawi pole puste, podstawiamy sensowną wartość domyślną „szt.”.

% =====================================================================
\section{\texttt{platformsdialog.cpp} --- okno zarządzania platformami}
% =====================================================================

Działa podobnie jak okno kategorii, więc pokazujemy tylko to, co tu jest inne ---
wybór typu nośnika z gotowej listy.

\subsection{Wybór z listy przez \texttt{QInputDialog::getItem}}

\begin{lstlisting}[caption={Gotowa lista rozwijana zamiast wpisywania}]
const QString typ = QInputDialog::getItem(
    this, "Typ nosnika", "Wybierz typ nosnika:",
    dostepneTypyNosnika(),   // {"Cyfrowy","Fizyczny","Streaming","Papier","E-book"}
    0,        // domyslnie zaznaczona pozycja
    false,    // false = lista zamknieta (uzytkownik nie wpisze swojego)
    &ok);
\end{lstlisting}

\textbf{Jak to działa:} \texttt{QInputDialog::getItem} pokazuje gotowe okienko
z listą rozwijaną. Podajemy zestaw dozwolonych typów nośnika, a parametr
\texttt{false} oznacza, że listy nie da się edytować --- użytkownik musi wybrać
jedną z gotowych opcji. To zabezpieczenie przed literówkami i bałaganem
(„Cyfrowy” vs „cyfrowe” vs „cyfra”).

\subsection{Domyślne zaznaczenie aktualnego typu}

\begin{lstlisting}[caption={Lista otwiera się na obecnej wartości}]
const QStringList typy = dostepneTypyNosnika();
// indexOf zwraca -1, gdy obecnego typu nie ma na liscie;
// qMax cofa wtedy do 0 (pierwsza pozycja), zeby nie podac blednego indeksu
const int indeksTypu = qMax(0, typy.indexOf(obecnyTyp));
const QString nowyTyp = QInputDialog::getItem(
    this, "Typ nosnika", "Wybierz typ nosnika:", typy, indeksTypu, false, &ok);
\end{lstlisting}

\textbf{Jak to działa:} przy edycji wygodnie jest, gdy lista od razu pokazuje
obecny typ platformy. \texttt{indexOf} znajduje pozycję obecnego typu na liście.
Jest jednak haczyk --- gdyby tego typu nie było na liście, \texttt{indexOf}
zwróciłby $-1$, co byłoby błędnym indeksem. Dlatego \texttt{qMax(0, ...)} pilnuje,
żeby wynik nigdy nie spadł poniżej zera --- w najgorszym razie zaznaczy się
pierwsza pozycja.

% =====================================================================
\section{\texttt{tagsdialog.cpp} --- okno zarządzania tagami}
% =====================================================================

Najprostsze z trzech okien zarządzania. Mimo to ma dwie warte uwagi rzeczy.

\subsection{Wygodne rozpakowanie pary (structured bindings)}

\begin{lstlisting}[caption={Rozpakowanie pary id--nazwa w pętli}]
const auto tagi = appController.pobierzTagi();   // lista par (id, nazwa)
int wiersz = 0;
for (const auto& [id, nazwa] : tagi) {           // rozpakowanie pary na 2 zmienne
    ui->tabela->insertRow(wiersz);
    ui->tabela->setItem(wiersz, 0, new QTableWidgetItem(QString::number(id)));
    ui->tabela->setItem(wiersz, 1, new QTableWidgetItem(nazwa));
    wiersz++;
}
\end{lstlisting}

\textbf{Jak to działa:} \texttt{pobierzTagi} zwraca listę par „(id, nazwa)”.
Zapis \texttt{for (const auto\& [id, nazwa] : tagi)} to tzw. \emph{structured
bindings} --- od razu rozkłada parę na dwie czytelne zmienne. Bez tego musielibyśmy
pisać \texttt{tag.first} i \texttt{tag.second}, co jest mniej zrozumiałe.

\subsection{Najprostsze usuwanie --- bez pytania o zawartość}

\begin{lstlisting}[caption={Usunięcie tagu kasuje tylko powiązania}]
// Tag nie pyta "co z zawartoscia" - usuwa tylko powiazania w tabeli
// posredniej multimedia_tagi. Same multimedia zostaja nietkniete.
const auto odp = QMessageBox::question(this, "Usun tag",
    QString("Czy na pewno usunac tag \"%1\"?\n"
            "Powiazania z mediami zostana usuniete.").arg(nazwa),
    QMessageBox::Yes | QMessageBox::No);
if (odp != QMessageBox::Yes) return;
appController.usunTag(idTagu);
\end{lstlisting}

\textbf{Jak to działa:} przy kategoriach i platformach trzeba było pytać, co zrobić
z przypisanymi pozycjami. Przy tagach nie ma tego problemu --- usunięcie tagu
kasuje tylko \emph{powiązania} w osobnej tabeli łączącej (\texttt{multimedia\_tagi}),
a same pozycje zostają nietknięte. Dlatego wystarczy zwykłe potwierdzenie Tak/Nie.

\textbf{Do czego służy:} pokazuje, że nie wszystko trzeba komplikować --- gdy dane
są luźno powiązane, prosta operacja w zupełności wystarcza.

% =====================================================================
\section{Podsumowanie}
% =====================================================================

Projekt CML, choć z pozoru to „zwykła aplikacja okienkowa”, kryje sporo ciekawych
rozwiązań. Po stronie bazy widać przemyślany SQL --- funkcje okna, transakcje
„wszystko albo nic” i dynamiczne sklejanie zapytań. Po stronie interfejsu
przewijają się powtarzalne, sprytne wzorce: chowanie danych w węzłach
(\texttt{Qt::UserRole}), automatyczne odświeżanie przez sygnały, czy
\texttt{blockSignals} ratujące przed nieskończonymi pętlami. Widać też dbałość
o drobiazgi (dymki, liczniki dni, kolory ostrzegawcze), które sprawiają, że
z programu po prostu wygodnie się korzysta.

\end{document}
